% ============================================================
% 00_packages.tex — Pakete und Paket-Level-Setup
% ============================================================
% Hinweis: Reihenfolge ist relevant. Schriften vor Mathe, xcolor mit
% [table]-Option vor allem, was Tabellen-Farben braucht.

\usepackage[T1]{fontenc}
\usepackage[utf8]{inputenc}
\usepackage[nswissgerman]{babel}
\usepackage[sfdefault,lf]{carlito}  % Calibri-Aequivalent (Google), lf = Lining Figures
\usepackage{sansmath}               % Mathe-Modus auch in Sans-Serif
\sansmath
\linespread{0.95}                   % Etwas engere Zeilen für Kompaktheit

\usepackage[landscape, a4paper, margin=1.2mm]{geometry} % MaschKo-Verdichtung: 1.5 -> 1.2mm (6-Seiten-Ziel)
\usepackage{microtype}
\usepackage{multicol}

\usepackage{mathtools}              % Superset von amsmath
\let\Bbbk\relax                     % Konfliktvermeidung mit newpxmath/amssymb
\usepackage{amssymb}
\usepackage{siunitx}
\sisetup{
  mode=match,
  propagate-math-font=true,
  reset-text-family=false,
  exponent-product = \cdot,
  output-decimal-marker = {,},
  per-mode = symbol
}

\usepackage[table]{xcolor}          % [table] aktiviert \rowcolor / \columncolor
\usepackage{transparent}
\usepackage{array}                  % erweiterte Spaltentypen + \arrayrulecolor
\usepackage{tabularx}               % \linewidth-Fit + X-Spalten
\usepackage{tabularray}             % moderne Tabellen (valuegrid)
\usepackage[column=O]{cellspace}    % cellspace nutzt 'O', damit siunitx 'S' frei bleibt
\usepackage{xstring}
\usepackage{varwidth}
\usepackage{ragged2e}               % Silbentrennung bei Flattersatz
\usepackage{enumitem}               % \setlist

\usepackage{graphicx}
\usepackage{tikz}
\usetikzlibrary{matrix,calc}
\usepackage{tcolorbox}
\tcbuselibrary{skins, breakable, raster}

\usepackage{hyperref}
